\heading{热学-课程说明}
\noteinfo{以下说明依据本项目现有真题、回忆稿与模拟题整理，主要反映题库中稳定出现的范围。}

\textbf{一、资料概况}

本项目当前收录本课程题目 9 套。整体看，题库覆盖了\textbf{热力学基本定律、理想气体与相变、气体动理论、统计分布、热力学判据与典型过程}。现有资料中既有选择 / 简答，也有一定比例的推导与计算题。

\textbf{二、考试范围（按现有题库归纳）}
\begin{itemize}[leftmargin=2em,itemsep=0.2em,topsep=0.2em]
  \item 热力学第一、第二定律，熵、自由能、焓、Gibbs 自由能等基本函数；
  \item Clausius 不等式、卡诺定理、可逆 / 不可逆过程比较；
  \item 理想气体、相变、节流过程、绝热过程与稳定性判据；
  \item Maxwell 速度分布、平均自由程、输运与动理论基础；
  \item Boltzmann 分布、重力场中的平衡分布与简单统计模型；
  \item 极端相对论气体、光子气体或与统计物理交叉的应用题。
\end{itemize}

\textbf{三、内容与命题特点}

热学题目的典型风格是：\textbf{基础概念题和推导题并存}。选择 / 简答题常考概念辨析，例如节流、相变、稳定性、热机效率；计算题则常要求你把热力学关系真正推到可用公式，而不仅仅是背结论。题目看似分散，但核心一直围绕“过程、状态函数与物理解释”三件事。

\textbf{四、难度评价}

整体难度可评为\textbf{中等}。如果公式之间关系不清，很容易觉得题目杂乱；但如果建立起“状态方程 -- 热力学势 -- 过程条件 -- 结论”的主线，很多题其实都能归到少数固定套路里。真正难的是\textbf{既要会算，又要会解释物理意义}。

\textbf{五、对课程的基本评价}

这门课很考验概念的内化程度。相比纯计算课程，热学更强调你是否真正理解熵、可逆性、平衡条件和状态函数之间的联系。现有题库的优点是代表性很强，拿来梳理知识框架尤其合适。

\textbf{六、如何更好利用模拟题}
\begin{itemize}[leftmargin=2em,itemsep=0.2em,topsep=0.2em]
  \item 把概念题和计算题分开刷：前者重在判断理由，后者重在推导路径；
  \item 每做完一道推导题，反过来问自己“这一步用了哪条基本关系”；
  \item 对节流、绝热、等温、相变等过程，建议单独整理比较表；
  \item 对动理论与统计分布部分，建议把 Maxwell 分布、Boltzmann 分布和几个常见平均量放在一页总结；
  \item 回忆稿与提纲稿中信息不完整的地方，不必死抠原题，应把它当作知识点提示来补全自己的复习网络。
\end{itemize}

